%% kapa: exam=1 round=28 year=2017 period=2 form=A subject=accounting subject_ko=회계학 seq=17 number=57
%% source: Q-Net 감정평가사 기출문제 (cst00309) artlSeq=5199603 file=extracted/1008/5199603/2017년 제28회 감정평가사 1차시험 문제지/2교시 A형.pdf
%% conversion: kordoc 4.24.51 -> tools/gen1.py (auto); review=proofread-sonnet; flags=none
\begin{question}{57}
유형자산의 감가상각에 관한 설명으로 옳지 않은 것은?
\begin{choices}
  \item 건물이 위치한 토지의 가치가 증가할 경우 건물의 감가상각대상금액이 증가한다.
  \item 유형자산을 수선하고 유지하는 활동을 하더라도 감가상각의 필요성이 부인되는 것은 아니다.
  \item 유형자산의 사용정도에 따라 감가상각을 하는 경우에는 생산활동이 이루어지지 않을 때 감가상각액을 인식하지 않을 수 있다.
  \item 유형자산의 잔존가치는 해당 자산의 장부금액과 같거나 큰 금액으로 증가할 수도 있다.
  \item 유형자산의 공정가치가 장부금액을 초과하더라도 잔존가치가 장부금액을 초과하지 않는 한 감가상각액을 계속 인식한다.
\end{choices}
\end{question}
\answer{1}
